\chapter{1986年上海卷（理）}

\section{单选题}

\begin{problem}

平移坐标轴，把原点 \(O(0,0)\) 移到 \(O'(2,-1)\)，则点 \((-1,-3)\) 在新坐标系中的坐标是（\quad）

\choices
    {\(\left( 3,2 \right)\)}
    {\(\left( -3,-2 \right)\)}
    {\(\left( -3,2 \right)\)}
    {\(\left( 3,-2 \right)\)}

\end{problem}

\begin{answer}
B
\end{answer}

\begin{problem}

下列三角关系式中不正确的是（\quad）

\choices
    {\(\sin\alpha+\sin\beta=2\sin \frac{\alpha+\beta}{2}\cos \frac{\beta-\alpha}{2}\)}
    {\(\sin\alpha-\sin\beta=2\cos \frac{\alpha+\beta}{2}\sin \frac{\beta-\alpha}{2}\)}
    {\(\cos\alpha+\cos\beta=2\cos \frac{\alpha+\beta}{2}\cos \frac{\beta-\alpha}{2}\)}
    {\(\cos\alpha-\cos\beta=2\sin \frac{\alpha+\beta}{2}\sin \frac{\beta-\alpha}{2}\)}

\end{problem}

\begin{answer}
B
\end{answer}

\begin{problem}

在直角坐标系中，直线 \(x+\sqrt{3}y+1=0\) 的倾角是（\quad）

\choices
    {\(\frac{\pi}{6}\)}
    {\(\frac{\pi}{3}\)}
    {\(\frac{2\pi}{3}\)}
    {\(\frac{5\pi}{6}\)}

\end{problem}

\begin{answer}
D
\end{answer}

\begin{problem}

函数 \(y=2\tan\left( 3x+\frac{\pi}{6} \right)\) 的最小正周期是（\quad）

\choices
    {\(\frac{\pi}{6}\)}
    {\(\frac{\pi}{3}\)}
    {\(\frac{\pi}{2}\)}
    {\(\frac{2\pi}{3}\)}

\end{problem}

\begin{answer}
B
\end{answer}

\begin{problem}

在空间，下列命题中正确的是（\quad）

\choices
    {如果两直线 \(a,b\) 与直线 \(l\) 所成的角相等，那么 \(a\parallel b\)}
    {如果两直线 \(a\)，\(b\) 与平面 \(\alpha\) 所成的角相等，那么 \(a\parallel b\)}
    {如果直线 \(l\) 与两平面 \(\alpha\)，\(\beta\) 所成的角都是直角，那么 \(\alpha\parallel\beta\)}
    {如果平面 \(\gamma\) 与两平面 \(\alpha\)，\(\beta\) 所成的二面角都是直二面角，那么 \(\alpha\parallel\beta\)}

\end{problem}

\begin{answer}
C
\end{answer}

\begin{problem}

当 \(|x|\leqslant 1\) 时，\(\arccos(-x)\) 等于（\quad）

\choices
    {\(\pi-\arccos x\)}
    {\(-\arccos x\)}
    {\(\arccos x\)}
    {\(\pi+\arccos x\)}

\end{problem}

\begin{answer}
A
\end{answer}

\begin{problem}

极坐标方程 \(\rho=\frac{5}{2-4\cos\theta}\) 所表示的曲线是（\quad）

\choices
    {椭圆}
    {抛物线}
    {双曲线}
    {直线}

\end{problem}

\begin{answer}
C
\end{answer}

\begin{problem}

下列双曲线方程中，以 \(y=\pm \frac{1}{2}x\) 为渐近线的是（\quad）

\choices
    {\(\frac{x^2}{16}-\frac{y^2}{4}=1\)}
    {\(\frac{x^2}{4}-\frac{y^2}{16}=1\)}
    {\(\frac{x^2}{2}-y^2=1\)}
    {\(x^2-\frac{y^2}{2}=1\)}

\end{problem}

\begin{answer}
A
\end{answer}

\begin{problem}

复数 \(\sqrt{3}-\i\) 的幅角的主值是（\quad）

\choices
    {\(\frac{\pi}{6}\)}
    {\(\frac{5\pi}{6}\)}
    {\(\frac{7\pi}{6}\)}
    {\(\frac{11\pi}{6}\)}

\end{problem}

\begin{answer}
D
\end{answer}

\begin{problem}

等式 \(\log_3 x^2=2\) 成立是等式 \(\log_3 x=1\) 成立的（\quad）

\choices
    {充分条件但不是必要条件}
    {必要条件但不是充分条件}
    {充分必要条件}
    {既不充分又不必要的条件}

\end{problem}

\begin{answer}
B
\end{answer}

\begin{problem}

在等比数列中，已知首项为 \(\frac{9}{8}\)，末项为 \(\frac{1}{3}\)，公比为 \(\frac{2}{3}\)，则项数是（\quad）

\choices
    {3}
    {4}
    {5}
    {6}

\end{problem}

\begin{answer}
B
\end{answer}

\begin{problem}

已知定直线 \(l\) 和外一定点 \(A\)，过点 \(A\) 且与 \(l\) 相切的圆的圆心轨迹是（\quad）

\choices
    {抛物线}
    {双曲线}
    {椭圆}
    {直线}

\end{problem}

\begin{answer}
A
\end{answer}

\begin{problem}

若 \(\sin \frac{\theta}{2}=\frac{3}{5}\)，\(\cos \frac{\theta}{2}=-\frac{4}{5}\)，则 \(\theta\) 角的终边在（\quad）

\choices
    {第一象限}
    {第二象限}
    {第三象限}
    {第四象限}

\end{problem}

\begin{answer}
D
\end{answer}

\begin{problem}

下列不等式中正确的是（\quad）

\choices
    {\(\sin \frac{5\pi}{7}>\sin \frac{4\pi}{7}\)}
    {\(\tan\frac{15\pi}{8}>\tan\left( -\frac{\pi}{7} \right)\)}
    {\(\sin \left( -\frac{\pi}{5} \right)>\sin \left( -\frac{\pi}{6} \right)\)}
    {\(\cos \left( -\frac{3\pi}{5} \right)>\cos \left( -\frac{9\pi}{4} \right)\)}

\end{problem}

\begin{answer}
B
\end{answer}

\begin{problem}

函数 \(y=\lg \left( x+\sqrt{x^2+1} \right)\)（\(x\in\R\)）（\quad）

\choices
    {是奇函数，不是偶函数}
    {是偶函数，不是奇函数}
    {既不是奇函数，又不是偶函数}
    {既是奇函数又是偶函数}

\end{problem}

\begin{answer}
A
\end{answer}

\begin{problem}

\(a\)，\(b\) 满足 \(0<a<b<1\). 下列不等式中正确的是（\quad）

\choices
    {\(a^a<a^b\)}
    {\(b^a<b^b\)}
    {\(a^a<b^a\)}
    {\(b^b<a^b\)}

\end{problem}

\begin{answer}
C
\end{answer}

\begin{problem}

正方体 \(ABCD-A_1B_1C_1D_1\) 中，直线 \(BC_1\) 与 \(AC\)（\quad）

\choices
    {相交且垂直}
    {相交但不垂直}
    {异面且垂直}
    {异面但不垂直}

\end{problem}

\begin{answer}
D
\end{answer}

\begin{problem}

设复数 \(z=a+b\i\)（\(a\)，\(b\in\R\)，且 \(b\ne 0\)），则 \(|z^2|\)，\(|z|^2\)，\(z^2\) 的关系是（\quad）

\choices
    {\(|z^2|=|z|^2\ne z^2\)}
    {\(|z^2|=|z|^2=z^2\)}
    {\(|z^2|\ne |z|^2=z^2\)}
    {互不相等}

\end{problem}

\begin{answer}
A
\end{answer}

\begin{problem}

过抛物线焦点 \(F\) 的直线与抛物线相交于 \(A\)，\(B\) 两点，若 \(A\)，\(B\) 在抛物线准线上的射影分别是 \(A_1\)，\(B_1\)，则 \(\angle A_1FB_1\) 等于（\quad）

\choices
    {\(45^\circ\)}
    {\(60^\circ\)}
    {\(90^\circ\)}
    {\(120^\circ\)}

\end{problem}

\begin{answer}
C
\end{answer}

\begin{problem}

若全集 \(I=\{(x,y)\mid x,y\in\R\}\)，\(A=\left\{(x,y)\mid \frac{y-3}{x-2}=1,\ x,y\in\R \right\}\)，\(B=\{(x,y)\mid y=x+1,\ x,y\in\R\}\)，则 \(\overline{A}\cap B\) 是（\quad）

\choices
    {\(\overline{A}\)}
    {\(B\)}
    {\(\varnothing\)}
    {\(\{(2,3)\}\)}

\end{problem}

\begin{answer}
D
\end{answer}

\section{填空题}

\begin{problem}

若三点 \(P_1\)，\(P_2\) 和 \(P\) 在一条直线上，点 \(P_1\) 和点 \(P_2\) 在直角坐标系中的坐标分别为 \((0,-6)\) 和 \((3,0)\)，且 \(\frac{P_1P}{PP_2}=-\frac{1}{2}\)，则点 \(P\) 的坐标是\fillinblank{}.

\end{problem}

\begin{answer}
\(\left( -3,-12 \right)\)
\end{answer}

\begin{problem}

正方体的全面积是 \(6\)，则其内切球体积是\fillinblank{}.

\end{problem}

\begin{answer}
\(\frac{\pi}{6}\)
\end{answer}

\begin{problem}

\(\lim\limits_{n\to\infty}\frac{1+3+5+\cdots+(2n-1)}{n(n+1)}=\)\fillinblank{}.

\end{problem}

\begin{answer}
\(1\)
\end{answer}

\begin{problem}

\(\left( \sqrt[3]{x}-\frac{1}{\sqrt{x}} \right)^8\) 的展开式中，\(x\) 的一次项的系数是\fillinblank{}.

\end{problem}

\begin{answer}
\(28\) 或 \(\mathrm C_8^6\)
\end{answer}

\begin{problem}

不等式组 \(\sqrt{x^2-9}\leqslant 4\) 且 \(x>0\) 的解是\fillinblank{}.

\end{problem}

\begin{answer}
\(3\leqslant x\leqslant 5\)
\end{answer}

\begin{problem}

轴截面是等边三角形的圆锥，它的侧面展开图扇形的圆心角的弧度数等于\fillinblank{}.

\end{problem}

\begin{answer}
\(\pi\)
\end{answer}

\begin{problem}

用 \(1\)，\(2\)，\(3\)，\(4\) 四个数字组成没有重复数字的四位奇数的个数是\fillinblank{}.

\end{problem}

\begin{answer}
\(12\)
\end{answer}

\begin{problem}

已知 \(\cos \theta=-\frac{3}{5}\) 且 \(\pi<\theta<\frac{3\pi}{2}\)，则 \(\cot\left( \theta-\frac{\pi}{4} \right)=\)\fillinblank{}.

\end{problem}

\begin{answer}
\(7\)
\end{answer}

\begin{problem}

函数 \(y=\log_2\left( \log_{\frac{1}{2}}x \right)\) 的定义域是\fillinblank{}.

\end{problem}

\begin{answer}
\((0,1)\)
\end{answer}

\begin{problem}

若直线的参数方程是 \(x=3+\frac{4}{5}t\)，\(y=-2+\frac{3}{5}t\)（\(t\) 为参数），则过点 \((4,-1)\) 且与其平行的直线在 \(y\) 轴上的截距是\fillinblank{}.

\end{problem}

\begin{answer}
\(-4\)
\end{answer}

\section{解答题}

\begin{problem}

设数列 \(\{a_n\}\) 的前 \(n\) 项的和为 \(S_n=n^2+2n+4\)（\(n\in\Z^+\)）.

\begin{enumerate}
\item 写出这个数列的前三项 \(a_1\)，\(a_2\)，\(a_3\)；
\item 证明：数列 \(\{a_n\}\) 除去首项后所成的数列 \(a_2\)，\(a_3\)，\(a_4\)，\(\cdots\) 是等差数列.
\end{enumerate}

\end{problem}

\begin{solution}
（1）由 \(a_1=S_1\)，得 \(a_1=1^2+2\times 1+4=7\).
又 \(a_2=S_2-S_1=2^2+2\times 2+4-7=5\)，\(a_3=S_3-S_2=3^2+2\times 3+4-\left( 2^2+2\times 2+4 \right)=7\).
所以 \(a_1=7\)，\(a_2=5\)，\(a_3=7\).

（2）当 \(n\geqslant 2\) 时，
\[
a_n=S_n-S_{n-1}=n^2+2n+4-\left[ \left( n-1 \right)^2+2\left( n-1 \right)+4 \right]=2n+1
\]
因此 \(a_{n+1}-a_n=\left[ 2\left( n+1 \right)+1 \right]-\left( 2n+1 \right)=2\).
故数列 \(a_2\)，\(a_3\)，\(a_4\)，\(\cdots\) 构成公差为 \(2\) 的等差数列.
\end{solution}

\begin{problem}

设 \(x>0\)，\(y>0\). 证明不等式 \(\left( x^2+y^2 \right)^{\frac{1}{2}}>\left( x^3+y^3 \right)^{\frac{1}{3}}\).

\end{problem}

\begin{solution}
只需证明 \(\left( x^2+y^2 \right)^3>\left( x^3+y^3 \right)^2\). 计算得
\begin{align*}
\left( x^2+y^2 \right)^3-\left( x^3+y^3 \right)^2
&=3x^4y^2+3x^2y^4-2x^3y^3\\
&=x^2y^2\left( 3x^2-2xy+3y^2 \right)\\
&=x^2y^2\left[ 3\left( x-\frac{y}{3} \right)^2+\frac{8}{3}y^2 \right]>0
\end{align*}
所以 \(\left( x^2+y^2 \right)^3>\left( x^3+y^3 \right)^2\). 由于两端都为正数，两端开 \(6\) 次方，得 \(\left( x^2+y^2 \right)^{\frac{1}{2}}>\left( x^3+y^3 \right)^{\frac{1}{3}}\).
\end{solution}

\begin{problem}

已知 \(k\) 为实数，复数 \(z=\cos\theta+\i\sin\theta\).

\begin{enumerate}
\item 当 \(k\) 和 \(\theta\) 分别为何值时，复数 \(z^3+k\overline{z}^3\) 是纯虚数；
\item 当 \(\theta\) 变化时，求出 \(\left| z^3+k\overline{z}^3 \right|\) 的最大值和最小值.
\end{enumerate}

\end{problem}

\begin{solution}
根据棣莫弗定理，
\begin{align*}
z^3+k\overline{z}^3
&=\cos 3\theta+\i\sin 3\theta+k\left( \cos 3\theta-\i\sin 3\theta \right)\\
&=\left( 1+k \right)\cos 3\theta+\i\left( 1-k \right)\sin 3\theta
\end{align*}

（1）若 \(z^3+k\overline{z}^3\) 是纯虚数，则
\[
\begin{cases}
\left( 1+k \right)\cos 3\theta&=0,\\
\left( 1-k \right)\sin 3\theta&\ne 0.
\end{cases}
\]
于是有两种情况：

\(k=-1\) 且 \(\theta\ne \frac{n\pi}{3}\)（\(n\in\Z\)）；

或 \(k\ne 1\) 且 \(\theta=\frac{2n\pi}{3}\pm \frac{\pi}{6}\)（\(n\in\Z\)）.

（2）由上式得
\begin{align*}
\left| z^3+k\overline{z}^3 \right|^2
&=\left( 1+k \right)^2\cos^2 3\theta+\left( 1-k \right)^2\sin^2 3\theta\\
&=1+k^2+2k\cos 6\theta
\end{align*}
因此：

当 \(k>0\) 时，\(\left| z^3+k\overline{z}^3 \right|\) 的最大值是 \(1+k\)，最小值是 \(|1-k|\)；

当 \(k=0\) 时，\(\left| z^3+k\overline{z}^3 \right|\) 的最大值和最小值都等于 \(1\)；

当 \(k<0\) 时，\(\left| z^3+k\overline{z}^3 \right|\) 的最大值是 \(1-k\)，最小值是 \(|1+k|\).
\end{solution}

\begin{problem}

如果抛物线 \(y^2=px\) 和圆 \((x-2)^2+y^2=3\) 相交，它们在 \(x\) 轴上方的交点为 \(A\) 和 \(B\)，那么当 \(p\) 为何值时，线段 \(AB\) 的中点 \(M\) 在直线 \(y=x\) 上.

\bitmapfigure[max width=0.42\linewidth]{img/1986/shanghai_science/q6_fig1.png}

\end{problem}

\begin{solution}
设 \(A(x_1,y_1)\)，\(B(x_2,y_2)\). 因为 \(A\)，\(B\) 同时在抛物线和圆上，所以由 \(y^2=px\) 代入 \((x-2)^2+y^2=3\)，得 \(x^2+\left( p-4 \right)x+1=0\).
故 \(x_1+x_2=4-p\)，\(x_1x_2=1\).

设 \(M(x_M,y_M)\)，则 \(x_M=\dfrac{x_1+x_2}{2}=\dfrac{4-p}{2}\).

又因为 \(A\)，\(B\) 在 \(x\) 轴上方，所以
\begin{align*}
y_1+y_2
&=\sqrt{\left( y_1+y_2 \right)^2}\\
&=\sqrt{y_1^2+y_2^2+2y_1y_2}\\
&=\sqrt{p(x_1+x_2)+2p\sqrt{x_1x_2}}\\
&=\sqrt{p(4-p)+2p}=\sqrt{6p-p^2}
\end{align*}
因此 \(y_M=\dfrac{y_1+y_2}{2}=\dfrac{\sqrt{6p-p^2}}{2}\).

因为 \(M\) 在直线 \(y=x\) 上，所以 \(\dfrac{\sqrt{6p-p^2}}{2}=\dfrac{4-p}{2}\).
化简得 \(6p-p^2=\left( 4-p \right)^2\)，即 \(p^2-7p+8=0\).
所以 \(p=\dfrac{7\pm \sqrt{17}}{2}\).
又因为 \(x_M=\dfrac{4-p}{2}\geqslant 0\)，故舍去 \(p=\dfrac{7+\sqrt{17}}{2}\). 因此 \(p=\dfrac{7-\sqrt{17}}{2}\).
\end{solution}

\begin{problem}

已知三角形 \(ABC\) 的两边 \(AC=2\)，\(BC=3\)，\(P\) 为边 \(AB\) 上一点（图1）. 现沿 \(CP\) 将此三角形折成直二面角 \(A-CP-B\)，当 \(AB=\sqrt{7}\) 时，求二面角 \(P-AC-B\) 的大小（图2）.

\bitmapfigure[max width=0.72\linewidth]{img/1986/shanghai_science/q7_fig1.png}

\end{problem}

\begin{solution}
在 \(\triangle ACB\) 中，\(BC=3\)，\(AC=2\)，\(AB=\sqrt{7}\). 因此 \(\cos \angle ACB=\dfrac{3^2+2^2-7}{2\times 3\times 2}=\dfrac{1}{2}\)，从而 \(\angle ACB=60^\circ\).

因为 \(A-CP-B\) 为直二面角，所以 \(\angle ACP+\angle PCB=90^\circ\)，且 \(\cos \angle ACP\cdot \cos \angle BCP=\cos \angle ACB=\dfrac{1}{2}\).
于是 \(\cos \angle ACP\cdot \cos \left( 90^\circ-\angle ACP \right)=\cos 60^\circ\)，即 \(\cos \angle ACP\cdot \sin \angle ACP=\dfrac{1}{2}\).
故 \(\sin 2\angle ACP=1\). 由 \(0^\circ<\angle ACP<90^\circ\)，得 \(\angle ACP=45^\circ=\angle BCP\).

自 \(P\) 作 \(PR\perp\) 平面 \(ACB\)，垂足为 \(R\). 因为 \(\angle ACP=\angle BCP\)，所以 \(R\) 在 \(\angle ACB\) 的平分线上.

自 \(R\) 作 \(RH\perp AC\)，连 \(PH\). 由三垂线定理可知 \(PH\perp AC\)，故 \(\angle PHR\) 为二面角 \(P-AC-B\) 的平面角.

\bitmapfigure[max width=0.46\linewidth]{img/1986/shanghai_science/q7_solution_fig2.png}

又因为 \(PH=CH\)，\(RH=CH\tan 30^\circ\)，所以
\[
\cos \angle PHR=\dfrac{RH}{PH}=\tan 30^\circ=\dfrac{\sqrt{3}}{3}
\]
因此二面角 \(P-AC-B\) 的大小为 \(\arccos \dfrac{\sqrt{3}}{3}\).
\end{solution}
